\section{A limitation of one-sided zeta substitution}\label{sec:substitution}
An upper rational approximation to $\zet$ increases a positive
$\zet$-coefficient matrix. It therefore solves the nonvanishing problem
without a perturbation or conditioning estimate. We now compare this
advantage with the necessary arithmetic cost just proved.

\subsection{A lower bound for the actual determinant}
Let $S=S_K$ and normalize the actual positive matrix as
\[
 T_K=C_S^{-1/2}M_K(\bet,\zet)_SC_S^{-1/2}>0.
\]
The bounds $4^n/(2n+1)\leq\binom{2n}{n}\leq4^n$ in
\eqref{eq:cpositive} give
\[
 \frac{7m}{256}\leq C_{mm}
 \leq\frac{7m(2K-2m+1)(2K+2m+1)}{256}.
\]
Hence $|\log\det C_S|=O(K\log K)$; it does not change a quadratic
logarithmic rate.

\begin{lemma}\label{lem:gramlower}
With
\begin{equation}\label{eq:Dstar}
 D_*=-\log2+\tfrac12\log(5/2)+\arctan(1/2)
                  +\tfrac14\arctan2+\tfrac\pi8,
\end{equation}
one has $\log\det T_K\geq-D_*K^2+o(K^2)$.
Here $D_*=0.8981320555\ldots$.
\end{lemma}
\begin{proof}
Only one fixed integration region is needed. If $t=\pi y$ and
$r=e^{-2t}$, direct differentiation gives
\[
 w(y)e^t=\frac{t^4}{3\pi}
 \frac{1+23r+23r^2+r^3}{(1-r)^4}\geq\frac1{48\pi},
\]
using $1-r\leq2t$. Also $J(4y^2)/O(4y^2)\geq1$. With $m_i=2i$,
\[
 (M_S)_{ij}=\frac1{16}\int_0^\infty
 \frac{J(4y^2)}{O(4y^2)}
 \frac{w(y)}{(y^2+m_i^2)(y^2+m_j^2)}\dd y.
\]
Expanding two determinants and applying Fubini (Andr\'eief's identity)
expresses $\det M_S$ as $16^{-d}/d!$ times an integral of the squared
Cauchy determinant and the $d$ weights. The squared determinant is
\begin{equation}\label{eq:cauchy-square}
 \frac{\prod_{i<j}(m_j^2-m_i^2)^2(y_j^2-y_i^2)^2}
      {\prod_{i,j}(m_i^2+y_j^2)^2}.
\end{equation}
Restrict the positive integral to the ordered boxes
\[
 \frac{sK(i-3/4)}d\leq y_i\leq\frac{sK(i-1/4)}d,
 \qquad 1\leq i\leq d,\qquad s=\tfrac12.
\]
Their volume and all scalar prefactors cost only $O(K\log K)$ in
the logarithm. The scaled pole and free coordinates approach uniform
distributions on $[0,1]$ and $[0,s]$ respectively, with mass
$\alpha=d/K\longrightarrow1/2$.

The logarithmic sums converge uniformly throughout these boxes.
Here are details at the singularities, where ordinary continuity
does not suffice. Scaled successive coordinates have separation at
least $c/d$, and the $i$th coordinate is at least $ci/d$.
A strip $|i-j|\leq\varepsilon d$ for the difference kernel contributes
at most
\[
 \frac{C}{d}\sum_{h\leq\varepsilon d}\log(Cd/h)
 =O\bigl(\varepsilon(1+|\log\varepsilon|)\bigr)+o(1)
\]
after division by $d^2$. At the origin use the lower bounds
$c(i+j)/d$ for a sum and $c_s(i^2+j^2)/d^2$ for the mixed squared
sum. The corner contribution is
$O(\varepsilon^2(1+|\log\varepsilon|))+o(1)$.
On the remaining compact regions the kernels are uniformly continuous.
First letting $d$ grow and then $\varepsilon$ decrease proves the
required convergence, including the missing Vandermonde diagonals.

The elementary integrals are
\begin{align*}
 I&=\int_{[0,1]^2}\log|x^2-v^2|\dd x\dd v=2\log2-3,\\
 J(s)&=\int_{[0,1]^2}\log(x^2+s^2v^2)\dd x\dd v\\
 &=\log(1+s^2)-3+\frac{\arctan s}{s}+s\arctan(1/s).
\end{align*}
For $I$, split into $\log|x-v|$ and $\log(x+v)$; for $J$, integrate
in each variable. Formula \eqref{eq:cauchy-square} and the weight
bound give
\[
 \liminf K^{-2}\log\det M_S
 \geq\alpha^2\{2I+2\log s-2J(s)\}-\frac{\pi\alpha s}{2}.
\]
The powers of $K$ in the products cancel up to $O(K\log K)$.
Substituting $\alpha=s=1/2$ gives $-D_*$. The normalization by
$\det C_S$ has lower-order logarithm, as observed above.
\end{proof}

\subsection{The certificate and its obstruction}
Suppose $0<p/q-\zet\leq q^{-2}$ and write $\epsilon=p/q-\zet$.
At the actual beta value the normalized matrix becomes
$T_K+\epsilon I$, which remains positive. Under the hypothetical
rationality $\bet=a/b$, the quantity
\[
 A_Kb^dq^dP_K(a/b,p/q)
\]
is a positive integer: each variable of $P_K$ has degree at most $d$.
Using only $\epsilon\leq q^{-2}$ gives the upper certificate
\begin{equation}\label{eq:envelope}
 \mathcal E_K(q;b)=A_Kb^d\det(qT_K+q^{-1}I).
\end{equation}
The global polynomial clearer $A_K$ is chosen before specializing $X$.
To obtain the proposed contradiction, this certificate would have to
tend to zero.

\begin{theorem}\label{thm:proxy}
For every fixed positive integer $b$,
\begin{equation}\label{eq:proxy-rate}
 \liminf_{K\to\infty}\frac1{K^2}
       \log\inf_{q>0}\mathcal E_K(q;b)
 \geq\gamma=\log2-D_*/2=0.2440811527\ldots>0.
\end{equation}
Thus the certificate \eqref{eq:envelope} cannot tend to zero for any
choice of positive $q=q_K$.
\end{theorem}
\begin{proof}
For every eigenvalue $\lambda_i>0$ of $T_K$,
$q\lambda_i+q^{-1}\geq2\sqrt{\lambda_i}$. Hence
\[
 \det(qT_K+q^{-1}I)\geq2^d\sqrt{\det T_K}
\]
uniformly in $q>0$. Combine this with \eqref{eq:necessaryA} and
Lemma~\ref{lem:gramlower}. The factor $b^d$ has logarithm $O_b(K)$.
The positive sign does not require a numerical certificate:
$\log2>2/3$, $\log(5/2)<1$, $\arctan(1/2)<1/2$,
$\arctan2<\pi/2$ and $\pi<22/7$ give $D_*<47/42$ and
$\gamma>3/28$.
\end{proof}

This conclusion concerns the displayed upper certificate. A lower
bound for an upper bound is not a lower bound for the actual integer
candidate. Exceptionally accurate approximations, cancellation after
specializing $X$ to a hypothetical rational value, and different
positive subspaces are not excluded. The theorem closes the specified
global-monic, standard-error argument; it does not close the
irrationality problem for $\bet$.
